% =============================================================================
% Methodology
% =============================================================================

\section{Methodology}
\label{sec:methodology}

\subsection{Uncertainty quantification}

Uncertainty quantification lies at the heart of our modeling approach. Disaster damage data typically exhibits positive skew and long right tails, where a small number of extreme events drive the majority of observed impacts. Consequently, traditional mean-based metrics fail to capture the reality of catastrophic risk. Instead, it is the tails of the distribution that demand attention. To address this, we employ generative probabilistic modeling \citep{gelman1995bayesian} as a framework for causal inference. This approach allows us to move beyond passive observation and perform counterfactual interventions using Judea Pearl's do-operator \citep{pearl2009causality}. By simulating the full distribution of potential outcomes under the intervention $P(\text{Damage} \mid \text{do}(\text{Policy}))$, we can conduct rigorous scenario analysis that accounts for functional relationships, prior research, and the inherent uncertainty in the data.

We specify probability distributions over model parameters (priors), update them with observed data via Bayes' theorem to obtain posterior distributions, and approximate these posteriors using Markov Chain Monte Carlo (MCMC) sampling---specifically the No-U Turn Sampler (NUTS) \citep{hoffman2014no}. Our central object of interest is the \textit{posterior predictive distribution}, which marginalizes over parameter uncertainty to produce a full probability distribution over possible outcomes. This allows us to quantify not just expected damages but also tail risks and extreme event probabilities, which are critical for building damage curves and their projections.

\subsection{Modeling disaster frequency}

As a first step, we model disaster frequency at country level. We use disaster data from the EM-DAT data set \citep{cred_emdat_2024}. We model the number of disasters a country experiences per year as a function of climate variables and country-level development indicators. We focus on one type of disasters: hydrometeorological (storms, floods, and movements of wet masses).

To model the number of disasters we use a Zero Inflated Binomial probability distribution, which allows us to better capture the fact that the generative process of disasters has a big mass concentration around zero, as multiple years do not report disasters for several countries. Consequently, we can say that the number of disasters occurring in country $i$ on year $t$ is distributed according to the Zero Inflated Binomial probability distribution:

\begin{equation}
Y_{i,t} \sim \text{ZeroInflatedBinomial}(\Psi, \mu, \alpha)
\end{equation}

The distribution is defined by three parameters: the inflation parameter $\Psi$ (probability of excess zeros), the mean $\mu$ (expected count conditional on the Negative Binomial process), and the dispersion $\alpha$. We model $\Psi$ using a Beta prior and $\alpha$ using a Gamma prior. The expected count $\mu = \EE[Y_{i,t} \mid X, \theta]$ is the core of the model, linked to observed covariates through a GLM structure. We use a softplus link function to ensure positive rates, computing $\mu_{i,t} = \text{softplus}(Z_{i,t})$ for every country $i$ in every year $t$. The latent rate $Z$ is defined as the sum of three components:

\begin{equation}
Z = \text{country\_fixed\_effect}_{i} + \text{spatial\_random\_effect}_{i} + \text{covariates\_effect}_{i,t}
\end{equation}

These components have the following meanings:

\begin{itemize}
    \item $\text{country\_fixed\_effect}$ is a model intercept with a country hierarchy. This means that it is modeled using a Bayesian hierarchical framework, which serves as an intermediary between the extremes of no-pooling (independent OLS per country) and complete pooling (a single global intercept). This structure assumes that country-level intercepts are drawn from a common global distribution, allowing for partial pooling of information across units. For countries with sparse data or extreme outliers, the hierarchical prior induces shrinkage of the country-specific estimates toward the global mean. This ``information sharing'' significantly improves the reliability of estimates for small-sample countries, effectively reducing the noise and over-fitting.

\item $\text{spatial\_random\_effect}_{i}$ accounts for spatial autocorrelation using a Besag-York-Molli\'{e} (BYM) specification \citep{besag1991bayesian}. The BYM decomposes the spatial effect into an unstructured term $\tilde{\theta}_i$ (capturing country-specific deviations) and a structured ICAR term $\phi_i$ (smoothing neighboring countries via an adjacency matrix $W$), combined as:
\end{itemize}

\begin{equation}
\text{spatial\_random\_effect}_{i} = \sqrt{1-\rho} \cdot \tilde{\theta}_i + \sqrt{\frac{\rho}{s}} \phi_i
\end{equation}

where $\rho$ is a mixing coefficient and $s$ a scaling factor for the ICAR variance. For the $\text{covariates\_effect}_{i,t}$ we introduce the time dimension $t$, and we use the following covariates:

\begin{itemize}
    \item $\text{ln\_population\_density\_standardized}_{i,t}$: log of population density for country $i$, period $t$.
    \item $\text{ln\_gdp\_pc\_standardized}_{i,t}$: log of GDP per capita for country $i$, period $t$.
    \item $\text{population\_standardized}_{i,t}$: Population in millions for country $i$, period $t$ squared.
    \item $\text{precip\_deviation\_standardized}_{i,t}$: Deviation of precipitation on country $i$ on period $t$ from its trend.
    \item $\text{co2\_standardized}_{t}$: Worldwide atmospheric CO2 concentration in parts per million at period $t$.
\end{itemize}

To control for unobserved heterogeneity on the panel structure, we apply the Mundlak adjustment by introducing group-level (country-level) random slopes for each covariate, which allows the model to account for correlation between covariates and latent group-specific effects. This way, the covariates functional form has the structure:

\begin{equation}
\text{covariates\_effect}_{i,t} = \sum_{j=0}^{n} X_{i,t,j} \cdot \beta_j + \sum_{j=0}^{n} \lambda_{i,j} \cdot \beta_j^{\text{latent}}
\end{equation}

Where:
\begin{itemize}
    \item $i$ corresponds to country index, $t$ year index, and $j$ covariate index.
    \item $\lambda_{i,j}$ is a country-specific deviation (random slope) for covariate.
\end{itemize}

\subsection{Modeling disaster damages}

For disaster damage, we use the logarithm of disaster damage. We model this variable using a Hurdle LogNormal probability distribution. The Hurdle LogNormal distribution is a type of hurdle model used to handle zero-inflated continuous data. Hurdle models are particularly useful when dealing with datasets that have an excess of zero values, in addition to continuous positive values. The Hurdle LogNormal model specifically combines a binary process that models the zeros and a truncated LogNormal process that models the positive continuous values.

The overall likelihood for the Hurdle LogNormal model is the product of the likelihoods from the binary component and the truncated LogNormal component:

\begin{itemize}
    \item $p$ if $y_i = 0$
    \item $(1-p) \cdot \text{LogNormal}(y_i|\mu, \sigma)$ if $y_i \neq 0$
\end{itemize}

Here, the LogNormal component is completely defined by the parameters $\mu$ and $\sigma$. To model sigma, we use a Half Normal distribution, which ensures strictly positive terms. While for $\mu$, we use exactly the same linear structure we used for frequency $\mu$ on equations (1) through (3).

\subsection{Modeling geospatial patterns of disasters}

As mentioned previously, the geospatial component is estimated using data from a specific region only, and predictions are generated only for a single target country at a time. This approach is necessary due to the substantial computational resources required to estimate the model at a global scale. To illustrate the methodology, we present results for Laos.

In the case of the probability of the event model, the objective is to model the probability of one event happening at any point of a country's grid. To achieve this, we create a set of synthetic non-disaster points, and we train the model to identify the probability of a point being a real disaster or a synthetic non-disaster point.

We estimate two different model versions. The first uses only the coordinates (latitude and longitude) of the grid and disaster points to capture geographical patterns of disaster occurrence. It captures these patterns by applying a Hilbert Space Gaussian Process (HSGP).

The second one incorporates, on top of the HSGP model, the distance to rivers and coastlines variables, plus a set of country-level covariates and climate-related time series.

\textbf{Hilbert Space Gaussian Process}

We approximate the spatial component using a Hilbert Space Gaussian Process (HSGP) \citep{solin2020hilbert}, which replaces the standard $O(N^3)$ Gaussian Process computation with a finite basis function expansion. The HSGP projects the GP onto a set of pre-computed basis functions derived from the spectral decomposition of the covariance kernel, reducing computational cost to $O(Nm)$ where $m \ll N$ is the number of basis functions. We use a Mat\'{e}rn-5/2 kernel, which produces smooth but flexible spatial surfaces, with learned length-scale and amplitude parameters. The HSGP takes latitude and longitude as inputs and learns the spatial pattern of disaster occurrence directly from the data, capturing regional variations without requiring explicit geographic covariates. This enables scalable inference over large spatial grids while preserving the uncertainty quantification properties of the full GP.

\textbf{Full model specification}

In this second step, we estimate the full model specification:

\begin{equation}
P(Y=1 \mid X)_{i,g,t} = \sigma\left( \text{country\_fix\_effect}_i + \text{HSGP}_{i,g} + X_{i,g,t} \cdot \beta_{i,g,t} \right)
\end{equation}

Here, in addition to the geospatial pattern detected by the HSGP, we obtain for every point in the grid of the region analyzed and for every disaster point the variables corresponding to 3 sets of covariates: climate change series, country-level specific variables, and distances to rivers and coastlines. All the covariates we introduce here are standardized to simplify the scale manipulation and the prior setting in the estimation. The standardization is done based on the whole model data set's mean and standard deviation. All the priors of the covariates parameters are defined as normal distributions with mean zero and standard deviation one.

The full set of covariates is:

\underline{Climate change series}:
\begin{itemize}
    \item $\text{co2\_standardized}_{t}$: Worldwide atmospheric CO2 concentration in parts per million at period $t$.
    \item $\text{dev\_ocean\_temp\_standardized}_t$: Deviation of the mean world's ocean temperature from its trend on period $t$.
    \item $\text{precip\_deviation\_standardized}_{i,t}$: Deviation of precipitation on country $i$ on period $t$ from its trend.
\end{itemize}

\underline{Country-level specific variables}:
\begin{itemize}
    \item $\text{ln\_population\_density\_standardized}_{i,t}$: log of population density for country $i$, period $t$.
    \item $\text{ln\_gdp\_pc\_standardized}_{i,t}$: log of GDP per capita for country $i$, period $t$.
    \item $\text{square\_ln\_gdp\_p\_standardized}_{i,t}$: log of GDP per capita for country $i$, period $t$ squared.
    \item $\text{population\_standardized}_{i,t}$: Population in millions for country $i$, period $t$ squared.
\end{itemize}

\underline{Distances to rivers and coastlines}:
\begin{itemize}
    \item For each point of the map, we compute the distance with the closest river, and we introduce its log as feature: $\text{log\_distance\_to\_river\_standardized}_{g}$.
    \item We also compute the distance to the closest coastline, and we include its log as feature: $\text{log\_distance\_to\_coastline\_standardized}_{g}$.
\end{itemize}

\subsection{Modeling patterns of climate and development indicators}

With the objective of being able to make predictions, we use a State Space modeling framework to predict trends in key inputs of the model. In particular, we use the implementation of State Space modeling via PyMC \citep{pymc-extras-statespace}. We project for a specific country of analysis: real GDP, population, deviation from trend precipitation. Additionally we project the global variables: atmospheric CO\textsubscript{2} accumulation, and ocean temperature deviations from trend.

For this, we employ a Bayesian dynamic linear model based on the PyMC state-space module, using a Bayesian Exponential Time Smoothing (ETS) \citep{ets_time_series} approach. This approach enables probabilistic inference and forecasting under a fully Bayesian framework. We model the multivariate time series using an additive trend exponential smoothing state-space model with the following configuration: local level (A), damped local trend (Ad), and no seasonal component (N). The model specification includes measurement error, a dense innovation covariance structure for the state shocks, and a stationary initialization with dampening. This setup provides flexibility to accommodate a wide range of time-varying behavior, including long-run stochastic trends, persistent shocks, and changing volatilities across the multiple indicators.

Formally, let $y_t \in \RR^n$ denote the observed vector of the $n$ variables at time $t$, which evolves according to a linear Gaussian state-space system:

\begin{align}
y_t &= Z\alpha_t + \epsilon_t, \quad \epsilon \sim \NN(0, H) \\
\alpha_t &= T\alpha_{t-1} + R\eta_t, \quad \eta_t \sim \NN(0, Q)
\end{align}

where $\alpha$ is the latent state vector (e.g., level and trend components), $H$ is the measurement error covariance matrix, and $Q$ is the innovation covariance matrix for the state shocks. These are learned from the data via a hierarchical prior specification, as described below.

\textbf{Prior Elicitation with Maximum Entropy Transformations}

To reflect informed but flexible beliefs about key smoothing and structural parameters, we employ maximum entropy--transformed priors (pz.maxent) for most hyperparameters. This allows us to specify constrained, information-efficient priors that respect known bounds or regularities (e.g., positivity, bounded damping).

\textbf{Multivariate Covariance via LKJ Prior}

We place an LKJ prior on the Cholesky decomposition of the state innovation covariance matrix $Q$, allowing for correlated innovations across state dimensions (e.g., co-movements between economic and climate variables). Specifically, we decompose $Q = LL^T$, where $L$ is the lower triangular Cholesky factor, and the standard deviations and correlations are separately inferred. We set the concentration parameter $\eta=4$, favoring weakly correlated structures but allowing for flexibility in learning strong correlations if supported by data. The standard deviation terms are assigned Gamma priors constrained to lie within plausible bounds.

\textbf{Forecasting}

Posterior Inference and Inverse Transformation: The model is then fit to data using Hamiltonian Monte Carlo (HMC) sampling via the nutpie backend \citep{nutpie2023} with jax acceleration \citep{jax2018}. The forecasts and trajectories generated by the model are returned in standardized space; hence, we apply an inverse exponential transformation composed with a StandardScaler inverse to convert model outputs back into the original data units.

\subsection{Country-level damage curves}

Once we computed the probability distributions of event frequency, and the probability distribution of damages from points 4.2 and 4.3, the expected damages can be easily computed as a product of both:

\begin{equation}
P(\text{expected\_damage})_{i,t} = P(\text{frequency})_{i,t} \times P(\text{damages})_{i,t}
\end{equation}

One advantage of this approach, and the Bayesian computations method we use, is that it allows us to propagate and quantify uncertainty from the initial two models to the expected damages.

Now that we have the expected damages, we resort to the notion of \textbf{return years} (RY) to represent the results in a much more intuitive way. A 100-RY disaster is one the associated damage of which is expected to be equaled or exceeded once every one-hundred years. Our primary object of study is, thus, a damage curve expressed in return years. The left panel of Figure~\ref{fig:example_curve} shows an example distribution over damages for a given year. The x-axis shows possible damage values, while the height is the probability density at that damage. Based on this curve, we can create a quantile function, which maps probabilities (on the x-axis) to the damage value which is larger than that percentage of values. The 100-RY damage value is larger than 99\% of all damage values produced by this distribution, and is marked.

\begin{figure}[H]
    \centering
    \caption[Example Damage Curve]{Example Damage Curve for a Single Year and Disaster Type.\\ \small The left curve shows all possible damage outputs for a category of disaster (e.g. flood, earthquake, drought, etc.). The x-axis shows the possible dollar-denominated damages, while the y-axis shows the probability density associated with that damage value (higher = more probable). The black bar shows the largest interval containing 95\% of all possible outcomes (called the 95\% HDI). The right curve shows the quantile of the distribution on the left. The 100-RY disaster, associated with the 99\% quantile, is marked.}
    \label{fig:example_curve}
    \includegraphics[width=\textwidth]{example_damage.png}
\end{figure}

Several features of Figure~\ref{fig:example_curve} are worth pointing out. The curve has a lot of mass around zero, followed by a long tail. Large disasters are, after all, relatively rare, with many years seeing none at all. As a result, the mean of the distribution (marked with an open circle on the left plot) is not at all interesting from a policy perspective. Instead, we are interested in quantifying the density in the \textit{tail}. Note that the 100-RY disaster, with a value of 4.1, is already a very rare event. But we can also contemplate even rarer events. In this example, a 500-RY event would have a value of 5.68. As the tail of the distribution becomes fatter, the difference between sizes of these rare events becomes larger.

Of primary interest to us is how climate change can potentially shift these curves, and how they vary according to geospatial patterns.

Climate change can shift these curves through two channels: disasters can become \textit{more frequent} (increasing the probability of occurrence $\mu$) or \textit{more severe} (shifting the damage distribution conditional on occurrence). The resulting damage distribution is a zero-inflated mixture: with probability $1 - \mu$ no disaster occurs and damage is zero; with probability $\mu$ a disaster occurs and damage follows the severity distribution. Expected damages are simply $\EE[D] = \mu \cdot \EE[D|I = 1]$, the product of frequency and average severity. Both channels increase total expected damages, and the causal framework presented in Figure~\ref{fig:causal_graph} justifies the statistical models we employ to estimate each component separately.

\subsection{Geolocated damage curves}

Once we have the country-level expected damages, we disaggregate these estimates spatially by leveraging our previously estimated probability of disaster occurrence at each geolocated grid point $g$. This allows us to construct geolocated damage curves, capturing how expected damages vary across space within each country. The key idea is to assign a share of national expected damages to each location, proportional to its relative disaster probability. This results in a spatially disaggregated risk map with expected damage curves for each grid point.

\begin{equation}
P(\text{expected\_damage})_{i,t,g} = \frac{P(Y=1|X)_{i,g,t}}{\sum_g P(Y=1|X)_{i,g,t}} \times P(\text{expected\_damage})_{i,t}
\end{equation}

This approach maintains internal consistency with the national-level expected damage figures. It distributes damages in proportion to the local probability of experiencing a disaster, as computed using our HSGP-based spatial event model. The denominator ensures the probabilities sum to one within each country-year, making the spatial distribution interpretable as a relative risk-weighted allocation. Importantly, this method inherits the full uncertainty structure of the Bayesian event and damage models, allowing us to construct local credible intervals around the expected damages at each location.

Once the expected damages are mapped at the local level, we proceed in a similar manner as before to construct geolocated damage curves. That is, for each location $g$, we generate the probability distribution over damage levels using the joint Bayesian posterior of both event probabilities and damage given an event. This yields a distribution $f_{D,g}(d)$ for each point $g$ from which we derive return-year quantiles (e.g., 100-RY or 500-RY events) for local areas.

Geolocated damage curves add critical granularity to the policy relevance of our framework. Different locations within a country face widely varying exposure and vulnerability profiles. For instance, coastal regions may face frequent flooding, while inland agricultural zones may suffer from drought. Our framework accounts for these spatial patterns by conditioning on local predictors such as distance to rivers, elevation, and population exposure, alongside the HSGP structure that learns spatial dependencies directly from the data.

Finally, a similar causal framework used to interpret country-level changes under climate scenarios applies locally. Increases in $P(Y=1|X)_{i,g,t}$ and $P(\text{events})_{i,t}$ derive in more frequent local disasters (e.g., more floods in coastal towns), while changes in $P(\text{damages})_{i,t}$ imply higher severity conditional on a disaster (e.g., stronger typhoons hitting a specific region). Spatially resolved return-year curves provide the basis for identifying which areas are most vulnerable, how risks change under counterfactual scenarios, and where targeted adaptation investments may be most needed.

\newpage
